%%%%%%%%%%%%%%%%%%%%%%%VICARIOUS%%%%%%%%%%%%%%%%%%%%%%%%%%%%%%%%%%%%%%%
% 																	%
% Template for presentation in Latex`s Beamer Class					%
% Using the default Berlin theme, can be replaced by other themes		%
% logo in the upper right can be replaced by new .png, gif, eps etc	%
% 																	%
%%%%%%%%%%%%%%%%%%%%%%%%%%%%%%%%%%%%%%%%%%%%%%%%%%%%%%%%%%%%%%%%%%%%%%%
\documentclass[xcolor=dvipsnames, aspectratio=169]{beamer}
\usetheme{Berlin}
\usecolortheme[named=LimeGreen]{structure}
\usepackage{beamerthemesplit} % kam neu dazu
\usepackage[ngerman]	{babel}			
\usepackage{t1enc}						
\usepackage[utf8]{inputenc}			
\usepackage{amsmath}
\usepackage{graphicx}
\graphicspath{{pictures/}}
\usepackage{amssymb}
\usepackage{amsfonts}
\usepackage{caption}
\usepackage{multimedia}
\usepackage{tikz}
\usepackage{listings}
\usepackage{acronym}
\usepackage{subfig}

\usepackage{lmodern}
\usepackage{multicol}

\definecolor{pblue}{rgb}{0.13,0.13,1}
\definecolor{pgreen}{rgb}{0,0.5,0}
\definecolor{pred}{rgb}{0.9,0,0}
\definecolor{pgrey}{rgb}{0.46,0.45,0.48}

\lstset{
    escapeinside={(*}{*)}
}


\lstdefinestyle{basic}{  
  basicstyle=\footnotesize\ttfamily,
  breaklines=true
  numbers=left,
  numberstyle=\tiny\color{gray}\ttfamily,
  numbersep=7pt,
  backgroundcolor=\color{white},
  showspaces=false,
  showstringspaces=false,
  showtabs=false,
  frame=single,
  rulecolor=\color{black},
  captionpos=b,
  keywordstyle=\color{blue}\bf,
  commentstyle=\color{gray},
  stringstyle=\color{green},
  keywordstyle={[2]\color{red}\bf},
}


\lstdefinelanguage{custom}
{
morekeywords={public, void},
sensitive=false,
morecomment=[l]{//},
morecomment=[s]{/*}{*/},
morestring=[b]",
}


\lstdefinestyle{BashInputStyle}{
  language=bash,
  showstringspaces=false,
  basicstyle=\small\sffamily,
  numbers=left,
  numberstyle=\tiny,
  numbersep=5pt,
  frame=trlb,
  columns=fullflexible,
  backgroundcolor=\color{gray!20},
  linewidth=0.9\linewidth,
  xleftmargin=0.1\linewidth
}

%Logo in the upper right just change if you know what you are doing^^
\addtobeamertemplate{frametitle}{}{%
\begin{tikzpicture}[remember picture,overlay]
\node[anchor=north east,yshift=2pt] at (current page.north east) {\includegraphics[height=1.8cm]{htw}};
\end{tikzpicture}}

\begin{document}
\bibliographystyle{alpha}
\title{Netzwerke -- Übung WiSe2018/19}
\subtitle{Application Layer\\
		\href{mailto:Benjamin.Troester@HTW-Berlin.de}{Benjamin.Troester@HTW-Berlin.de}\\
		PGP: ADE1 3997 3D5D B25D 3F8F 0A51 A03A 3A24 978D D673 }
\author{Benjamin Tröster}

\date{}

\begin{frame}
\titlepage

\end{frame}

\section*{Road-Map}
\begin{frame}
\frametitle{Road-Map}
\begin{multicols}{2}
  \tableofcontents
\end{multicols}
\end{frame}

\section{Retrospektive}
\begin{frame}{Retrospektive}
\begin{itemize}[<+->]
	\item Vorlesung
	\begin{itemize}
		\item Retrospektive der Vorlesung -- was haben Sie behandelt?
		\item Fragen?
	\end{itemize}
\end{itemize}
\end{frame}


\section{Präsentation}
\subsection{DNS}
\begin{frame}
	\frametitle{1.) DNS}
	\begin{itemize}
		\item Erläutern Sie kurz was DNS ist und welche Komponenten das DNS im wesentlichen benötigt um Domainnamen aufzulösen.
		\item Wie können Domainnamen aufgelöst werden? Wie sieht die Auflösung in der Realität aus?
		\item Erläutern Sie anhand eines Beispiels, wie die Namensauflösung stattfinden kann.
		\item Erläutern Sie die Semantik und Syntax der Tool \emph{dig} und \emph{host}. 
		\item Wozu kann (konnte) das Tool \emph{whois} genutzt werden?
	\end{itemize}
\end{frame}

\section{Krypto}
\begin{frame}
	\frametitle{2.) Krypto I}
	\begin{itemize}
		\item Erläutern Sie kurz, was unter den Begriffen Kryptologie, Kryptografie, Kryptoanalyse verstanden wird.
		\item Welche Aufgaben sollen mithilfe Kryptografie gewährleistet werden?
		\item Welche Maßnahmen setzen die eben genannten Aufgaben um?
		\item Erläutern Sie was unter dem Begriff Kerckhoff'sche Prinzip verstanden wird.
	\end{itemize}
\end{frame}

\begin{frame}
	\frametitle{3.) Krypto II}
	\begin{itemize}
		\item Erläutern Sie kurz welche beiden Verschlüsselungsverfahren ganz grundsätzlich unterschieden werden.
		\item Was wird unter dem Begriff Hashing verstanden, wie wird diese Idee in der Kryptografie angewandt?
		\item Erläutern Sie die Aufgabe von kryptografischen Signaturen!
	\end{itemize}
\end{frame}	
	
\section{SSH}
\begin{frame}
	\frametitle{4.) SSH}
	\begin{itemize}
		\item Erläutern Sie wo die grundlegenden Konfigurationsdateien liegen und was deren Aufgabe ist.
		\item Was ist ein Fingerprint in Bezug auf SSH?
		\item Lassen sich die kryptografischen Schlüssel die in SSH erzeugt wurden wiederverwenden? Oder muss für jeden Teilnehmer ein neues Schlüsselpaar generiert werden?
		\item Erläutern Sie was unter einem VPN verstanden wird!
	\end{itemize}
\end{frame}	

\end{document}
